\section{Introduction}

Topological complexity, introduced by Farber \cite{Farber2003}, measures the discontinuity of motion planning.
For a path-connected space $X$, it is the (unreduced) sectional
category of the endpoint fibration
\[
 \pi:PX\longrightarrow X\times X,\qquad
 \pi(\gamma)=(\gamma(0),\gamma(1)).
\]
For Riemannian manifolds, geodesic complexity imposes the additional
requirement that the chosen paths be minimizing geodesics
\cite{RecioMitter2021,MescherStegemeyer2023}.

Let $(M,g)$ be a closed negatively curved manifold of dimension
$n\geq 2$. The universal covering space $\widetilde M$ is a Hadamard
manifold, so fixed-endpoint homotopy classes determine unique geodesics.
Different classes can nevertheless contain minimizing geodesics of
equal length. Write $m_g(x,y)$ for their number.

We study how the multiplicity of minimizing geodesics reflects the
topology of $M$. Our primary goal is a quantitative lower bound,
involving simplicial volume, for the number of pairs admitting the
maximal number of minimizing geodesics under a general-position
condition. The topological complexity of closed negatively curved
manifolds is already known \cite{Dranishnikov2020}; our aim is not to reprove its value, but
to refine the geometric information associated with motion planning.

\subsection{Generic multiplicity}

Let $\mathcal R_-(M)$ denote the space of smooth negatively curved
metrics, with the $C^\infty$ topology, and put $\Gamma=\pi_1(M)$.
For lifts $u,v\in\widetilde M$ of $x,y\in M$ and $a\in\Gamma$, set
\[
 L_a(u,v)=d_{\widetilde M}(u,av),\qquad
 S_g(u,v)=\{a\in\Gamma:L_a(u,v)=d_M(x,y)\}.
\]
The set of minimizing labels $S_g(u,v)$ is finite, and
$m_g(x,y)=|S_g(u,v)|$.

\begin{definition}
A negatively curved metric $g$ is called \emph{regular} if, whenever
$S_g(u,v)=\{a_0,\ldots,a_{k-1}\}$ has $k\geq2$ elements, the map
\[
 (L_{a_1}-L_{a_0},\ldots,L_{a_{k-1}}-L_{a_0})
\]
has differential of rank $k-1$ at $(u,v)$.
\end{definition}

For a regular metric, $m_g(x,y)\leq2n+1$. Let
\[
 C_g=\{(x,y)\in M\times M:m_g(x,y)=2n+1\}.
\]
We regard these as ordered pairs.

\begin{theorem}\label{thm:intro-count}
Let $M$ be a connected, closed, oriented smooth $n$-manifold,
$n\geq2$, admitting a negatively curved metric.
Regular metrics form an open dense subset of $\mathcal R_-(M)$.
For each regular metric $g$, the set $C_g$ is finite and
\[
 \#C_g\geq c_n\|M\|^2,\qquad
 c_n=\begin{cases}
 \displaystyle\frac{2}{n^2((n+1)!)^2},&n\text{ even},\\[5pt]
 \displaystyle\frac{1}{2n(2n+1)!},&n\text{ odd}.
 \end{cases}
\]
Here $\|M\|$ denotes simplicial volume in the sense of
Gromov \cite{Gromov1982}.
\end{theorem}

\subsection{Exterior canonical classes and localization}

Put $G=\Gamma\times\Gamma$ and let $P=\Z[\Gamma]$ carry the action
$(g,h)\cdot a=gah^{-1}$. The augmentation kernel
$I=\ker(P\to\Z)$ supports the Costa--Farber canonical class
$\mathfrak v\in H^1(G;I)$ \cite{CostaFarber2010}. If $q_k:I^{\otimes k}\to\bigwedge^k I$
is the exterior quotient, define
\[
 \mathfrak a_k=q_{k*}(\mathfrak v^k),\qquad
 A_M=\mathfrak a_{2n}\frown[M\times M]
 \in(\bigwedge^{2n}I)_G.
\]
On $V_k=(\bigwedge^k(I\otimes\R))_G$, we use the quotient
$\ell^1$-seminorm $\|\cdot\|_\Delta$ whose generators of cost one are
\[
 [(a_1-a_0)\wedge\cdots\wedge(a_k-a_0)].
\]

\begin{theorem}\label{thm:intro-estimate}
For $M$ as in Theorem~\ref{thm:intro-count},
\[
 c_n\|M\|^2\leq\|A_M\|_\Delta\leq\|M\times M\|.
\]
\end{theorem}

The central geometric construction is the nerve of the sets of
simultaneously minimizing geodesic labels. For a regular metric its
top-dimensional simplices correspond to the pairs in $C_g$.
We obtain a localization formula expressing $A_M$ as a signed sum
of the associated exterior augmentation generators. Consequently
$\|A_M\|_\Delta\leq\#C_g$, which combines with
Theorem~\ref{thm:intro-estimate} to give
Theorem~\ref{thm:intro-count}.

These local terms have a metric-homological interpretation.
We use magnitude homology of metric spaces as developed in
\cite{LeinsterShulman2021}; its relationship with interval
order complexes is studied in \cite{KanetaYoshinaga2021}.
For distinct $x,y\in M$, the endpoint-fixed magnitude homology
at the minimal length satisfies
\[
 \MH_{2,d_M(x,y)}(M;x,y)\cong
 \ker(\Z[S_g(u,v)]\longrightarrow\Z).
\]
In particular, its rank equals $m_g(x,y)-1$; at pairs in $C_g$
the top exterior power is generated by the corresponding local
term in the localization formula.

\subsection{Bounded cohomology and geodesic complexity}

The estimate for $\|A_M\|_\Delta$ uses bounded cohomology and
its relation to simplicial volume \cite{Gromov1982}.
For the $G$-set $\Gamma\cong G/\Delta\Gamma$, the common stabilizer
of two distinct labels is a cyclic group: it is isomorphic to the
centralizer of a nontrivial element of $\Gamma$. This amenability
property permits controlled lifts of certain bounded classes
from $G$ to bounded invariant alternating cochains on $\Gamma$.

For regular metrics, smooth multiplicity strata also yield the
geodesic-complexity bound
\[
 \GC(M,g)\leq 2n+1.
\]
Together with the known value $\TC(M)=2n+1$
(in our unreduced convention) \cite{Dranishnikov2020}, this gives
\[
 \GC(M,g)=\TC(M)=2n+1
 =1+\max_{x\ne y}\operatorname{rank}
 \MH_{2,d_M(x,y)}(M;x,y).
\]
Here $\GC$ is defined using the unreduced convention.

The remainder of the paper is organized as follows. Section 2 studies canonical classes and magnitude homology.
Section 3 introduces the minimizing-geodesic nerve.
Section 4 develops the bounded cohomological lifting argument,
and Section 5 proves quantitative seminorm estimates.
Section 6 treats generic metrics and geodesic complexity.
Section 7 establishes localization and the counting theorem.
Section 8 discusses further questions.
